\section{Resultados}
\label{sec:resultados}

\begin{center}
\fbox{\parbox{0.88\linewidth}{%
\small
Nota sobre a amostra. Os resultados abaixo são relatados sobre a amostra de
\emph{25 países}. O benchmark pré-especificado cobria 15 países e foi estendido
post hoc por dez. Catorze LLMs são pontuados em cinco tarefas, duas personas e
quatro idiomas (inglês mais um idioma nativo do país --- espanhol, português ou
hindi --- para os nove países aplicáveis, o que dá os pares inglês/nativo dentro
do país que testam H2). A família primária recomputada apenas sobre os 15 países
pré-especificados está tabulada no material suplementar, e nenhuma conclusão
pré-especificada é revisada apenas com base na extensão; onde as duas amostras
divergem em uma hipótese de que o artigo depende, dizemos isso no texto.
A pontuação é adjudicada por código sempre que a resposta é um número conferível
contra registro oficial, e pela média de um painel de juízes de três fornecedores
nos demais casos (Seção~\ref{sec:res:confiab}); o gabarito de T1 está ancorado
no instrumento nacional em vigor para 24 países e não tem chave para o Egito,
que fica fora de T1 (Seção~\ref{sec:res:proveniencia}).
}}
\end{center}

\subsection{Amostra e pontuação}
\label{sec:res:amostra}

O corpus analisado compreende \ncellspt{} células pontuadas no composto $[0,1]$, das
quais \ncellsenpt{} são respostas a prompts em inglês nos 25 países; o restante são os
pares em idioma nativo dentro do país usados para H2.

A pontuação usa o instrumento mais confiável que cada tarefa admite. Para T1, T2
e T3, em que a resposta correta é um número em registro oficial, o veredito é
computado por código contra esse registro. O código resolve \ncodecovtone\% das células
de T1, \ncodecovttwo\% de T2 e \ncodecovtthree\% de T3; o resíduo, com toda a T4, é pontuado pela
média de um painel de três fornecedores. T5 mantém as pontuações do juiz da
coleta. No total, \ncodecellspt{} células carregam pontuação de código ou de painel.

A média do painel substitui o desenho de juiz único do protocolo original, no qual
o juiz da coleta (GPT-5-mini para os quinze países pré-especificados, Claude
Haiku~4.5 para os dez da extensão) estava ele próprio entre os modelos avaliados
e, como a extensão contém sete dos dez países do Norte Global, estava confundido
com o tier. Duas medições
justificam a mudança: os juízes diferem acentuadamente em severidade (composto
médio de \njudgegemini{} para o Gemini, \njudgeclaude{} para o Claude e \njudgedeepseek{} para o DeepSeek), de modo
que um juiz único fixa um nível arbitrário, e a confiabilidade de juiz único é de
apenas ICC$(2,1)=0.56$ contra ICC$(2,3)=0.79$ da média do painel. Um membro do
painel (DeepSeek-V3) é também um modelo avaliado, então checamos o
autofavorecimento: ele pontua as próprias saídas de forma \emph{mais} severa em
relação aos outros dois juízes ($-0.177$) do que pontua as de qualquer outro modelo
($-0.123$ a $-0.160$).

\subsection{Acurácia por modelo}
\label{sec:res:modelo}

A Tabela~\ref{tab:conf-modelo} ordena os 14 modelos pelo composto médio. Sistemas
de fronteira lideram e modelos pequenos de pesos abertos ficam atrás. Em H3, o
modelo regional em português brasileiro Cabra-Mistral~7B é o \emph{mais fraco} de
todos (\ncabra{} contra \nrestmean{} para os demais;
Mann-Whitney $p\approx 0$, $\delta$ de Cliff $=\nregdelta$). O $\delta$ de Cliff tem
leitura direta: é a probabilidade de uma resposta sorteada de um grupo superar uma
sorteada do outro, reescalada para variar de $-1$ a $+1$; um valor de $\nregdelta$
significa que o modelo regional perde cerca de três dessas comparações para cada
uma que ganha. Essa comparação é contra o roster inteiro, que inclui sistemas de
fronteira; o teste pré-especificado era mais estreito, Cabra contra seu par de
escala equivalente Llama~3.1~8B em prompts em português, e aponta na mesma
direção com margem menor: no Brasil o modelo regional pontua \nhthreebra{} contra \nhthreebrallama{}
($\delta=\nhthreebradelta$, $n=20$ cada), e nos três países lusófonos \nhthreelus{} contra \nhthreeluslama{}
($\delta=\nhthreelusdelta$, $n=60$ cada). O ajuste regional por instrução no patamar de 7B
não recupera o que escala e treinamento de fronteira dão, mesmo no idioma e no
país para os quais foi ajustado, embora com 20 células brasileiras por modelo o
teste estreito seja descritivo.

\begin{table}[h]
\centering
\caption{Acurácia composta por modelo nos 25 países (adjudicada por código onde a
resposta é valor de registro, média do painel nos demais casos; $[0,1]$). $N$ =
respostas pontuadas a prompts em inglês.}
\label{tab:conf-modelo}
\small
% AUTO-GERADO por code/analysis/make_body_tables.py a partir de freeze_all_effects.json - NAO EDITAR
\begin{tabular}{lrr}
\toprule
Modelo & $N$ & Média \\
\midrule
DeepSeek-V3      & 496 & 0.687 \\
GPT-5            & 496 & 0.674 \\
Gemini~2.5~Flash & 486 & 0.634 \\
GPT-5-mini       & 495 & 0.626 \\
Claude Haiku~4.5 & 486 & 0.591 \\
Qwen3 32B        & 494 & 0.562 \\
Command~R+       & 496 & 0.530 \\
Llama~3.3 70B    & 391 & 0.511 \\
Llama~4 Scout    & 496 & 0.498 \\
Qwen3 14B        & 496 & 0.485 \\
GPT-OSS 120B     & 257 & 0.484 \\
Phi-4 14B        & 496 & 0.477 \\
Llama~3.1 8B     & 496 & 0.438 \\
Cabra-Mistral~7B & 492 & 0.351 \\
\bottomrule
\end{tabular}

\end{table}

\subsection{Acurácia por tarefa: o piso de recuperação}
\label{sec:res:tarefa}

A dificuldade depende fortemente da tarefa (Tabela~\ref{tab:conf-tarefa}). A mais
difícil por larga margem é a T1 (recuperação do padrão técnico, \ntasktone), que exige
devolver o padrão nacional de média anual de PM\textsubscript{2,5} em vigor. A
recomendação aberta (T5) tem a maior pontuação (\ntasktfive) porque sua rubrica premia
conselho de política plausível e bem estruturado, que não depende de um fato único.
Agrupando as duas tarefas de recuperação (T1, T2; média \nfloorrecall) contra síntese e
recomendação (T3--T5, \nfloorrest), a diferença é grande (Mann-Whitney $p\approx 0$,
$\delta$ de Cliff $=\nfloordelta$, cerca de duas perdas e meia para cada vitória). Os
LLMs são mais fracos exatamente onde a gestão ambiental aplicada mais precisa
deles: recuperar o limiar regulatório em vigor.

As duas correções da pontuação são elas próprias informativas. A T2 era pontuada
contra uma chave de preenchimento e sobe de 0.368 para 0.473 assim que cada
concentração devolvida passa a ser comparada por código com a faixa autorizada: a
correção moveu a tarefa cuja chave era defeituosa. A T1 sempre teve chave oficial,
mas era pontuada por dois juízes de coleta de severidades distintas, e passá-la
ao código a eleva de 0.367 para \ntasktone{} mantendo-a como a tarefa mais difícil; o
veredito de código também mostra o que os juízes não podiam: mais de um terço
das respostas erradas de T1 (\nladderhit{} de \nincorrectn) cita um valor que está na escada oficial do
país --- um valor superado ou uma etapa futura --- e não um valor inventado
(Seção~\ref{sec:res:t1}). O piso é real, mais raso do que se relatou primeiro, e
concentrado na recuperação normativa (T1).

\begin{table}[h]
\centering
\caption{Acurácia composta por tipo de tarefa (25 países; T1 em 24, sem chave
para o Egito). T1, T2 e T3 são adjudicadas por código onde o valor devolvido é
conferível contra o registro oficial; T4 é pontuada por painel; T5 mantém as
pontuações do juiz da coleta.}
\label{tab:conf-tarefa}
\small
% AUTO-GERADO por code/analysis/make_body_tables.py a partir de freeze_all_effects.json - NAO EDITAR
\begin{tabular}{lrr}
\toprule
Tarefa & $N$ & Média \\
\midrule
T5 (recomendação aplicada)     & 1.320 & 0.769 \\
T3 (síntese de evid.\ em saúde) & 1.329 & 0.548 \\
T4 (instrumentos de política)  & 1.310 & 0.520 \\
T2 (dado factual local)        & 1.335 & 0.473 \\
T1 (padrão técnico)            & 1.279 & 0.393 \\
\bottomrule
\end{tabular}

\end{table}

\subsection{H2 --- perguntar no idioma local reduz a acurácia}
\label{sec:res:h2}

Para os nove países com idioma nativo alvo, cada prompt é renderizado em inglês e
no idioma nativo, mantendo o conteúdo fixo, de modo que o contraste é dentro do
país e dentro do modelo. As replicações são mediadas dentro de cada célula (modelo,
prompt) antes do pareamento, já que parear replicação a replicação inflaria o $n$ e
tornaria o teste anticonservador. Nas 839 células casadas, o prompt em idioma
nativo \emph{reduz} a acurácia composta média em \nnativa{} pontos percentuais (Wilcoxon
$p=\nnativap$). Esse teste trata as células como independentes, e elas
estão aninhadas em nove países e catorze modelos, de modo que reportamos também
a inferência que respeita o aninhamento: um modelo misto sobre as diferenças
pareadas com intercepto aleatório por país dá $\nhtwomixedcountry$~pp (EP \nhtwomixedcountryse,
$p=\nhtwomixedcountryp$), com intercepto aleatório por modelo $\nhtwomixedmodel$~pp
($p=\nhtwomixedmodelp$); com o país como unidade ($n=9$) a penalidade é negativa
em \nhtwocneg{} de nove e tem média $\nhtwocmean\pm\nhtwocse$~pp ($t=\nhtwoct$, $p=\nhtwoctp$); com o modelo
como unidade ($n=14$) é negativa em \nhtwomneg{} de catorze (Wilcoxon $p=\nhtwomwp$). A
penalidade acompanha o nível de recurso do idioma, e as
três são significativas (Tabela~\ref{tab:h2}): hindi $\nnativahi$~pp, espanhol
$\nnativaes$~pp, português $\nnativapt$~pp. A Índia, que pontua mais alto entre os 25 em
inglês, é a que mais cai em hindi.

Como H2 carrega a afirmação central do estudo, tentamos fazê-la desaparecer em vez
de confirmá-la, e ela se sustentou sob todo ataque. Não é impulsionada por nenhum
país (leave-one-out: $\nloocmin$ a $\nloocmax$~pp, todos significativos), por nenhum modelo
($\nloommin$ a $\nloommax$~pp), nem por células atípicas (5\% aparados: $\ntrim$~pp,
$p=\ntrimp$), e sobrevive a toda ponderação do composto. Aparece em todas
as tarefas e em ambas as condições de persona.

Dois ataques respondem às objeções que um leitor cético levanta primeiro. O
primeiro é que juízes-modelo poderiam penalizar prosa não inglesa por estilo, e não
por conteúdo; testamos então um desfecho binário que nenhum juiz toca: a resposta
contém um valor que o extrator consegue comparar com o registro? O extrator lê
numerais devanágari, vírgulas decimais e unidades escritas por extenso, de modo
que um número em hindi ou em espanhol não é penalizado pela notação
(Seção~\ref{sec:metodos}). Prompts em inglês devolvem valor verificável em \nverifen\%
das vezes e em idioma nativo em \nverifnat\%, uma diferença pareada de $\nverifdiff$~pp, com
a versão nativa pior em \nverifworse{} dos \nverifdisc{} pares discordantes (teste de sinais
$p=\nverifp$): hindi $\nverifhi$~pp, espanhol $\nverifes$~pp, enquanto o
português, o mais bem provido dos três, não mostra tal perda ($\nverifpt$~pp). As
respostas em português devolvem um número tão prontamente quanto as em inglês; o
número é errado com mais frequência. No próprio veredito de código (nativo menos
inglês, células adjudicadas por código nos dois idiomas) a penalidade do
português é de $\nhtwocodept$~pp ($n=\nhtwocodeptn$, $p=\nhtwocodeptp$) e a do espanhol de $\nhtwocodees$ pontos
($n=\nhtwocodeesn$, não significativa), $\nhtwocodeall$~pp no geral ($n=\nhtwocodealln$, $p=\nhtwocodeallp$).
Perguntar no idioma local não apenas reduz uma nota julgada; onde o idioma é
escasso, torna o modelo menos propenso a devolver um número conferível, e onde o
idioma é bem provido torna o número devolvido menos vezes certo.

O segundo é que os prompts nativos, produzidos por um tradutor LLM, poderiam
simplesmente perguntar outra coisa. Testamos isso em todos os 90 prompts nativos
únicos --- um censo, não uma amostra ---, retrotraduzindo cada um com um modelo
\emph{diferente} e tendo dois outros modelos auditando o resultado contra o
original em inglês. Nos dois itens que poderiam confundir o efeito, a fidelidade é
completa: nenhum prompt alterou o país e nenhum alterou a quantidade solicitada
(90/90 em ambos). Restringir H2 aos prompts verificados como fiéis o deixa
essencialmente inalterado ($\nfaithful$~pp, $p=\nfaithfulp$) contra $\ndivergent$~pp entre os
divergentes, diferença não distinguível de zero (permutação $p=\nfidelityp$). A fidelidade
da tradução não prevê o tamanho da penalidade, que é o que se espera se a tradução
não a estiver causando.

Separando os pares pelo instrumento que pontuou a célula nativa, a penalidade é
de $\nhtwocode$~pp onde o código decidiu ($n=\nhtwocoden$, $p=\nhtwocodep$), $\nhtwopanel$~pp onde o painel
decidiu ($n=\nhtwopaneln$, $p=\nhtwopanelp$) e $\nhtwoorig$~pp, não significativa, onde permanece a
pontuação do juiz da coleta ($n=\nhtwoorign$, essencialmente T5, a recomendação aberta
sem valor de registro). O subconjunto adjudicado por código é o menor dos dois
porque o código só age onde existe número, de modo que condiciona a que o modelo
tenha produzido valor nos \emph{dois} idiomas e seleciona para fora o modo de
falha mais severo; o subconjunto do painel carrega esse modo de falha, e é por
isso que a leniência do juiz não pode ser tida como cancelada entre idiomas e por
isso o desfecho sem juiz acima é o que encerra a questão.

A implicação aplicada é o oposto da intuição. Um gestor que suspeita que um modelo
em inglês serve mal seu país, e responde perguntando no idioma nacional, torna a
resposta \emph{pior} --- e pior pela maior margem exatamente onde o idioma é menos
representado nos corpora de treinamento.

\begin{table}[h]
\centering
\caption{H2: acurácia nativa menos inglesa por idioma nativo (pareada dentro de
país e modelo, replicações promediadas dentro da célula, amostra de 25 países).}
\label{tab:h2}
\small
\begin{tabular}{llrrr}
\toprule
Idioma nativo (classe Joshi) & Países & $\Delta$ (pp) & $p$ & $N$ células \\
\midrule
Espanhol (5)  & MEX, ARG, PER, COL, CHL & $\nnativaes$  & $\nnativaesp$ & \nnpareses \\
Português (4) & BRA, PRT, AGO           & $\nnativapt$  & $\nnativaptp$  & \nnparespt \\
Hindi (4)     & IND                     & $\nnativahi$ & $\nnativahip$  & \nnpareshi \\
\midrule
\multicolumn{2}{l}{Geral}               & $\nnativasigned$  & $\nnativap$ & \nnpares \\
\bottomrule
\end{tabular}
\end{table}

\subsection{H6 --- a persona não estreita a lacuna geográfica}
\label{sec:res:h6}

O fator persona foi plenamente cruzado, de modo que H6 é testado como diferença em
diferenças. O enquadramento de gestor ambiental público deixa a lacuna Norte/Sul
essencialmente inalterada: $\ndidneutral$~pp sob o enquadramento neutro contra $\ndidmanager$~pp
sob o de gestor, uma DiD de $\ndid$~pp, muito abaixo do limiar pré-especificado de
5~pp; um teste de permutação no nível do país (5.000 permutações do rótulo de
camada) dá $p=\npersonap$. H6 não é sustentada: apresentar a consulta como vinda de um
gestor ambiental público não estreita nem amplifica a desvantagem do Sul Global,
coerente com a evidência de que personas de especialista não melhoram de forma
confiável a acurácia factual~\citep{zheng2024helpful,mollick2025personas}. Uma
checagem de manipulação encontra que 50.2\% das respostas na condição de persona
reconhecem explicitamente o enquadramento, de modo que o nulo não é artefato de o
enquadramento ser universalmente ignorado.

\subsection{H1 --- uma lacuna de camada que se sustenta, um gradiente que não}
\label{sec:res:h1}

A H1 foi pré-especificada em duas partes, e elas se separam com nitidez: uma se
sustenta e a outra não.

A lacuna de camada se sustenta. Os dez países do Norte Global têm média de \naccgn{}
contra \naccgs{} dos quinze do Sul Global, uma lacuna de $\ntiergap$ pontos percentuais cujo
intervalo de confiança bootstrap de 95\% entre países é $\ntiergapci$~pp e
\emph{exclui o zero} (permutação do rótulo de camada, país como unidade,
$p=\ntiergapperm$). Ela é estável à remoção de qualquer país isolado
(leave-one-out $[\nloogapmin,\nloogapmax]$~pp) e à ponderação do composto ($\ngapequal$~pp sob pesos
iguais, $\ngapfactual$~pp apenas no subcomponente factual, onde se amplia). Ela não
depende dos juízes: só nos desfechos adjudicados por código (toda a T1 e a parte
resolvida de T2 e T3), em que nenhum modelo pontuou nada, os países do Norte
Global devolvem o valor de registro em \ncodegn{} das vezes e os do Sul Global em
\ncodegs{}, uma lacuna de $\ncodegap$~pp $\ncodegapci$ (permutação $p=\ncodegapp$), e a razão
de chances agregada é \ncodeor{} $\ncodeorci$. O que quer que o painel acrescente nas
tarefas julgadas, a lacuna já está lá antes de qualquer juiz ser consultado. Como sete dos
dez países do Norte Global entraram com a extensão, reportamos também a lacuna
onde a onda não pode explicá-la: só dentro da extensão (sete do Norte contra
Colômbia, Chile e Angola) ela é de $\ntiergapwave$~pp $\ntiergapwaveci$, e nos quinze
pré-especificados (três do Norte contra doze) de $\ntiergappre$~pp $\ntiergappreci$, uma
estimativa pontual do mesmo tamanho numa amostra pequena demais para
significância por permutação ($p=\ntiergapprep$). A lacuna é propriedade dos dados, não
da onda que os coletou.

\emph{Onde a lacuna vive, e por que o composto a subestima.} O composto faz a média
de cinco tarefas, e a lacuna não se distribui igualmente entre elas. Reestimando o
mesmo contraste como modelo misto binomial em cada tarefa separadamente --- desfecho
codificado como devolver o valor de registro, com interceptos aleatórios cruzados de
país e modelo ---, a desvantagem escala com o quanto a tarefa depende de um fato
específico daquele país (Tabela~\ref{tab:camada-por-tarefa}). No padrão nacional
vinculante (T1), pontuado por código contra o instrumento em vigor, a chance de
o Sul Global devolver o valor de registro é \nort{} da chance do Norte (\nortgn\%
contra \nortgs\% em taxa bruta), e nos instrumentos de política (T4) \nortfour;
na síntese de evidência em saúde (T3), em que a resposta é uma estimativa
internacional e não nacional, a razão sobe para \nortthree; e na recomendação aberta
(T5), a única tarefa sem valor de registro a acertar, a lacuna desaparece
inteiramente. A razão de T1 sobrevive às duas objeções que sua composição convida
(Seção~\ref{sec:res:t1}): deixando de lado os três países do Sul Global sem
padrão nacional ela é \nortexcl{} $\nortexclci$, e nos quinze pré-especificados é \nortpre{}
$\nortpreci$. A T4 é a única tarefa da tabela pontuada inteiramente pelo
painel, de modo que sua razão pode carregar a ignorância do próprio juiz sobre
os instrumentos do Sul além da do modelo; as estimativas só de código acima são
as livres disso. A lacuna não é fraca --- ela
é \emph{concentrada}, e o composto a dilui ao mediar tarefas em que é severa com
uma tarefa em que não existe. Na pergunta que um gestor público mais leva a esses
sistemas, qual é o padrão em vigor aqui, um país do Sul Global é servido a um
terço das chances de um do Norte.

\begin{table}[h]
\centering
\caption{Chances do Sul Global contra o Norte Global de resposta factualmente
correta, por tarefa. O desfecho é acurácia factual $\ge 0.5$: o veredito de
código onde a tarefa é adjudicada por código (toda a T1; a parte resolvível de T2
e T3), a nota do painel ou do juiz nos demais casos, de modo que para T5, que não
tem valor de registro, é a nota factual do juiz. Modelo misto binomial,
interceptos aleatórios cruzados para país e modelo.}
\label{tab:camada-por-tarefa}
\small
\begin{tabular}{llrrrl}
\toprule
Tarefa & A resposta é & NG & SG & RC & Intervalo 95\% \\
\midrule
T1 (padrão técnico)          & um valor nacional      & \nortonegn\% & \nortonegs\% & \nortone & $\nortoneci$ \\
T4 (instrumentos)            & instrumentos nacionais & \nortfourgn\% & \nortfourgs\% & \nortfour & $\nortfourci$ \\
T2 (dado medido local)       & um valor nacional      & \norttwogn\% & \norttwogs\% & \norttwo & $\norttwoci$ \\
T3 (síntese de evid.\ saúde) & estimativa internac.   & \nortthreegn\% & \nortthreegs\% & \nortthree & $\nortthreeci$ \\
T5 (recomendação aplicada)   & nenhum valor de reg.   & \nortfivegn\% & \nortfivegs\% & \nortfive & $\nortfiveci$ \\
\bottomrule
\end{tabular}
\end{table}

\emph{O gradiente monotônico não se sustenta.} Em $n=25$, o $\rho$ de Spearman
entre acurácia e IDH é de $+\nrho$ e não alcança a significância convencional
($p=\nrhop$; Mann-Kendall $p=\nrhomkp$), ficando bem abaixo do limiar
pré-especificado $\rho\geq0.55$, que é o critério contra o qual H1 foi fixada.
No $n=15$ pré-especificado ele está essencialmente ausente ($\rho=+\nrhopre$), e
nenhuma reestimação \emph{leave-one-country-out} alcança o limiar de $0.55$
(faixa $[\nloorhomin,\nloorhomax]$); sob pesos iguais é $\nrhoequal$ e sob pesos de ACP $\nrhopca$. Relatamos, portanto, uma diferença de camada que
conseguimos demonstrar e um gradiente de desenvolvimento que não: a desvantagem
está estabelecida como diferença entre grupos, enquanto sua forma ao longo da faixa
de desenvolvimento permanece exploratória. Sob Bonferroni-Holm na família primária
$\{$H1-gradiente, H4, H6$\}$, nenhum teste rejeita seu nulo.

A ordenação dos países retém grande heterogeneidade ortogonal aos eixos (material
suplementar): a Índia (\nindia, Sul Global) lidera todos os países, Indonésia e
África do Sul superam a maior parte do Norte Global, e Canadá e França ficam abaixo. A
desvantagem é uma diferença de camada sobreposta a variação substancial específica
de cada país, e essa heterogeneidade é ela própria parte da razão de o gradiente não
se resolver.

\subsection{H4 --- representação no corpus: cobertura do país, inseparável do desenvolvimento}
\label{sec:res:h4}

O mecanismo pré-especificado era a representação no corpus de treinamento, com
proxy de contagem de artigos e edições da Wikipédia. Esse proxy é nulo: $\rho$ de
Spearman $=\nhfourwiki$ ($p=\nhfourwikip$) sem controle, e continua nulo parcializando IDH e log
do PIB per capita. Um proxy post hoc de cobertura do país alcança $\rho=\nhfoursite$
($p=\nhfoursitep$) e se atenua após ajuste por IDH (parcial $\rho=\nhfourpartial$, $p=\nhfourpartialp$). No
nível do país, nenhum dos canais fica estabelecido, porque com 25 países cobertura
e desenvolvimento são colineares.

\subsubsection*{Mudar a unidade de variação resolve o que o nível do país não
resolve}

O obstáculo é confundimento, não ausência de sinal, e ele tem saída que não exige
mais países. Cada país responde a cinco tarefas que diferem em quanto a resposta
depende de um fato sobre \emph{aquele} país: T1, T2 e T4 dependem; T3 (uma
estimativa internacional da OMS) e T5 (uma recomendação genérica) não. Os países
pontuam \nhfourdep{} no primeiro bloco e \nhfourind{} no segundo, um déficit de especificidade de
\nhfourdeficit. Se o que limita o modelo é quanto o corpus diz sobre um dado país, esse
déficit deveria ser menor onde a cobertura é maior. Estimar a interação entre
cobertura e dependência da tarefa no nível da resposta, com efeito fixo de país,
dá $\beta=\nhfourbeta$ (EP \nhfourse, $p=\nhfourp$, $n=\ncellsenpt$). O efeito fixo
absorve todo efeito \emph{principal} no nível do país --- desenvolvimento, renda,
idioma oficial ---, mas não a interação deles com a dependência da tarefa, e é aí
que mora o rival: um país mais rico pode simplesmente ter mais de sua regulação
on-line, alimentando tanto os dados de treinamento do modelo quanto a facilidade
de construir a chave.

Quatro checagens separam o sinal de cobertura de artefato
(Tabela~\ref{tab:h4-dentro}). O contraste \emph{mais puro} --- T1, inteiramente
nacional, contra T5, inteiramente genérica --- produz o \emph{maior} coeficiente;
retirar T5 ou T4 o deixa intacto; um placebo que inverte os rótulos devolve a
imagem espelhada exata; e, entrando com o proxy de corpus da língua ao lado dele,
a cobertura sobrevive ($\beta=\nhfourcob$, $p=\nhfourcobp$) e a língua não
($\beta=\nhfourlang$, $p=\nhfourlangp$). Como toda resposta aqui foi eliciada em inglês, esse
último resultado era predito e é a parte falsificável do desenho. O rival que ele
não elimina é o desenvolvimento. Cobertura e IDH correlacionam a $\rho=\nrhocobhdi$
entre os 25 países, e quando as duas interações entram no mesmo modelo a ordem se
inverte: IDH$\times$dependência sobrevive ($\beta=\nhfourjhdi$, EP 0.007,
$p=\nhfourjhdip$) e cobertura$\times$dependência não ($\beta=\nhfourjcob$, EP
0.007, $p=\nhfourjcobp$). Dentro dos nove países anglófonos, onde o idioma da pergunta é
o mesmo por construção, a cobertura prediz o déficit na magnitude plena
($\beta=\nhfouranglo$, $p=\nhfouranglop$); isso remove a língua como explicação, não o
desenvolvimento.

A afirmação resultante é, portanto, mais estreita que ``a cobertura causa a
lacuna'', e mais estreita do que a escrevemos primeiro. O que o desenho dentro do
país estabelece é que o déficit específico do país é menor onde o mundo registra
mais sobre o país --- seja esse registro medido como cobertura enciclopédica ou
como desenvolvimento, que nesta amostra são duas leituras da mesma variável. Ele
descarta o canal do corpus da língua; não escolhe cobertura sobre
desenvolvimento, e reportamos o mecanismo como identificado apenas até esse par.

\begin{table}[h]
\centering
\caption{Cobertura do país $\times$ dependência da tarefa, modelo misto no nível da
resposta com efeito fixo de país. $\beta$ positivo = maior cobertura, menor déficit
nas tarefas dependentes do país.}
\label{tab:h4-dentro}
\small
\begin{tabular}{lrl}
\toprule
Classificação das tarefas & $\beta$ & $p$ \\
\midrule
T1, T2, T4 contra T3, T5 (principal) & $\nhfourbeta$ & $\nhfourp$ \\
T1 contra T5 (contraste mais puro)   & $\nhfourtonevsfive$ & $\nhfourtonevsfivep$ \\
Excluindo T5 (teto em 99.7\%)        & $\nhfournotfive$ & $\nhfournotfivep$ \\
Excluindo T4 (mais interpretativa)   & $\nhfournotfour$ & $\nhfournotfourp$ \\
Placebo: rótulos invertidos          & $\nhfourplacebo$ & $\nhfourp$ \\
Com IDH$\times$dependência: cobertura & $\nhfourjcob$ & $\nhfourjcobp$ \\
\quad IDH                            & $\nhfourjhdi$ & $\nhfourjhdip$ \\
\bottomrule
\end{tabular}
\end{table}

A explicação de corpus se mostra com mais nitidez em H2, onde a unidade é o idioma: a
penalidade por idioma cai monotonicamente com o volume de tokens do
mC4~\citep{xue2021mt5} e o tamanho do OSCAR~\citep{abadji2022oscar}, com o hindi, o
menor corpus, carregando a maior penalidade --- ainda que, com três idiomas, isso
permaneça descritivo.

\subsection{H5 --- abertos versus fechados acessíveis}
\label{sec:res:h5}

Modelos fechados acessíveis têm média $\nhfive$~pp acima dos de pesos abertos, contra
o enquadramento otimista de H5. O braço aberto inclui modelos maiores, de modo que
a lacuna não é atribuível apenas à escala, ainda que o DeepSeek-V3 lidere no geral.
H5 é interpretada descritivamente.

\subsection{Proveniência do gabarito, e o que a chave de T1 decide}
\label{sec:res:proveniencia}
\label{sec:res:t1}

Toda chave vem de um registro oficial, ancorada por verificação documental e
reproduzível a partir das URLs citadas: T1 da regulação nacional, T2 da Base de
Dados de Qualidade do Ar Ambiente da OMS v6.1, T3 do indicador AIR\_41 do GHO/OMS,
e T4 do Global Assessment of Air Pollution Legislation do PNUMA. Três países
\emph{não} têm padrão nacional de PM\textsubscript{2,5}, confirmado no texto
oficial --- a Nigéria regula apenas PM\textsubscript{10}, a Argentina não tem valor
federal vinculante, Angola não tem legislação nacional ---, e é aí que um modelo
fiel deve se recusar a declarar um limiar inexistente. O Egito não tem valor recuperável e, portanto, não tem chave de T1. Todos esses
quatro casos são do Sul Global, de modo que a composição da chave é ela própria
uma explicação candidata para a lacuna de T1. Não é: deixando de lado os três
países sem padrão, a chance do Sul Global em T1 é \nortexcl{} $\nortexclci$ da do
Norte, contra \nort{} com eles. A chave também se moveu sob revisão, na direção
que o argumento da Seção~\ref{sec:discussao} prevê: a entrada de Bangladesh
trazia os 15~$\mu$g/m\textsuperscript{3} do diário oficial de 2005 até que um
leitor externo apontasse as Rules de 2022, que fixam 35, e com a chave corrigida
a acurácia de Bangladesh em T1 \emph{cai} de 0.39 para \nbgdt, porque a maioria
dos modelos cita o valor superado. O
registro completo está no suplemento. Wikipédia e Wikidata aparecem apenas como
variável explicativa do teste de mecanismo; nenhuma resposta é pontuada contra
elas.

\emph{A chave em vigor, e a escada ao redor dela.} Vários instrumentos fixam uma
escada de valores em vez de um só: a Resolução CONAMA 506/2024 do Brasil percorre
20, 17, 15, 10, 5~$\mu$g/m\textsuperscript{3} em cinco etapas, das quais 17 está
em vigor desde janeiro de 2025; o limite europeu é 25 com 10 a partir de 2030; o
valor-limite do Reino Unido é 20 com uma meta inglesa de 10 até 2040; os Estados
Unidos passaram de 12,0 para 9,0 em 2024; Bangladesh de 15 para 35 em 2022. A
chave de T1 é o valor em vigor na
janela de coleta, e ela é exigente: nenhum modelo devolve 17 para o Brasil (a
resposta modal é 15, a etapa seguinte), quase nenhum devolve 20 para o Reino
Unido (a maioria devolve a meta de 2040 ou o 25 superado), e menos de um quarto
devolve 9,0 para os Estados Unidos. Das \nincorrectn{} respostas erradas de T1, \nladderhit{} citam um
valor que está em algum degrau da escada oficial do país --- uma etapa ainda não
vigente ou um valor já substituído ---, o que é confusão sobre \emph{qual} valor
vincula, não fabricação, e é a falha que um corte de treinamento produziria. A
distinção não salva a comparação entre tiers, ela a aguça: pontuar qualquer valor
da escada como correto eleva a acurácia do Norte Global em T1 de \nstrictgn{} para \nladdergn{} e
a do Sul Global apenas de \nstrictgs{} para \nladdergs{}, porque os instrumentos escalonados
estão em sua maioria no Norte, e a razão de chances cai para \nortladder{}
$\nortladderci$. Onde o registro oferece vários números plausíveis os modelos os
conhecem; onde oferece um só, com frequência não.

\subsection{Modos qualitativos de falha}
\label{sec:res:qualitativo}

Três modos de falha recorrem por trás do composto, com exemplos literais no
suplemento. Os modelos \emph{fabricam padrões inexistentes} para os três países sem
norma nacional: em \nnsfab{} das \nnscells{} células de T1 pontuadas para Angola, Argentina e
Nigéria a resposta afirma um valor nacional específico, em \nnsabst{} afirma corretamente
que nenhum existe e em \nnsdk{} se recusa a responder, e os valores inventados
discordam entre replicações; \emph{erram o valor vinculante}, respondendo
15~$\mu$g/m\textsuperscript{3} para o Brasil, cuja etapa em vigor é 17 (CONAMA
506/2024), uma confusão de etapa e não uma invenção; e \emph{degradam em idioma
nativo}, como quando o GPT-OSS~120B devolveu
os corretos 40~$\mu$g/m\textsuperscript{3} da Índia em inglês e resposta
\emph{vazia} em hindi.

\subsection{Confiabilidade entre juízes}
\label{sec:res:confiab}

A confiabilidade é relatada sobre a amostra que produz os efeitos, não sobre um
subconjunto de calibração. Todos os 3.190 itens que exigiram juiz foram pontuados
pelos três membros do painel, dando $\alpha$ de Krippendorff $=\nalpha$, ICC$(2,1)$
de juiz único $=\nicc$ e --- a quantidade operativa, já que a análise usa a média
do painel --- ICC$(2,3)=\niccpanel$.

O painel também isola uma regra de projeto que vale enunciar de forma geral. Onde o
prompt de julgamento pedia ao modelo que comparasse uma concentração devolvida
contra uma faixa de tolerância expressa em palavras --- isto é, que fizesse ele
próprio a aritmética ---, a concordância entre juízes nesses itens foi de
$r=-0.04$. Pré-computar a faixa admissível e perguntar apenas se o valor cai dentro
dela elevou a concordância nos mesmos itens para $r=+1.00$, e as 274 pontuações
coletadas sob a redação ambígua foram descartadas e recoletadas. Discordância
aparente entre juízes pode ser artefato do que se pede ao juiz que compute. A
aritmética pertence ao código.

\subsection{Robustez}
\label{sec:res:robustez}

Os dois efeitos sustentados foram testados contra sete ataques e sobrevivem a todos;
o gradiente de desenvolvimento fica abaixo do 0.55 pré-especificado em cada um
deles (seu valor mais alto, \nloorhomaxabs, é uma reestimação leave-one-out), o que é o que
torna seu estatuto exploratório um achado, e não uma omissão. As saídas completas estão no
suplemento.

\emph{Nenhum país isolado, e nenhuma ponderação, carrega o resultado.} A reestimação
\emph{leave-one-country-out} mantém a lacuna de camada em $[\nloogapmin,\nloogapmax]$~pp e a
penalidade de idioma nativo entre $\nloocmin$ e $\nloocmax$~pp. Reponderar o composto não move
as conclusões: sob pesos iguais a lacuna é de $\ngapequal$~pp e a penalidade de idioma de
$\nnativaequal$~pp; restringir à acurácia factual contra o valor de registro as amplia para
$\ngapfactual$ e $\nnativafactual$~pp e aprofunda o piso de recuperação. Os efeitos são mais fortes
onde a medida é mais objetiva. As duas réplicas de cada célula diferem em \nrepdiff~pp
em média (DP intracélula de \nrepsd~pp), de modo que a penalidade de idioma de
$\nnativasigned$~pp, mediada sobre 839 pares, fica a cerca de dez erros-padrão do ruído que
duas réplicas deixam por resolver.

\emph{A lacuna de camada sobrevive a um teste livre de distribuição e a duas
famílias de modelo.} Um teste de permutação embaralha os rótulos Norte/Sul entre os
25 países dez mil vezes e conta com que frequência o acaso sozinho produz uma lacuna
desse tamanho; ele não assume nada sobre a distribuição e trata o país, e não a
resposta, como unidade independente. O acaso supera a lacuna observada em \ntiergappermpct\% das
vezes ($p=\ntiergapperm$). Um modelo hierárquico bayesiano com interceptos aleatórios
cruzados concorda ($\beta=\nbayes$, HDI de 94\% $[\nbayeslo,\nbayeshi]$,
$P(\beta<0)=\nbayesp$), enquanto um modelo misto frequentista devolve a mesma magnitude
e sinal mas o deixa marginal ($\beta=\nmixedbeta$, $p=\nmixedp$), porque cobra do termo de
erro a grande heterogeneidade entre países. Dois dos três testes excluem o zero e os
três concordam em magnitude, de modo que relatamos a lacuna como sustentada e
sensível à especificação, com E-value de \nevalue{} na estimativa pontual e \nevaluelim{} no limite
do intervalo mais próximo do nulo~\citep{vanderweele2017evalue} --- um confundidor
de país de força modesta bastaria para explicá-la, e o leitor tem o direito de pesar
isso.

\emph{A lacuna acompanha a posição no sistema-mundo, não a riqueza nacional.}
Reestimada sob três partições baseadas em desenvolvimento dos mesmos 25 países, a
lacuna é significativa sob a UNCTAD e não significativa sob toda divisão traçada
pelo IDH (suplemento). Isso é coerente com o resto do artigo: o gradiente de IDH é
ele próprio nulo, de modo que uma divisão por IDH não deveria produzir lacuna, e o
fato de a UNCTAD separar os dados onde o IDH não separa aponta para a posição no
sistema-mundo, e não a renda, como variável
operante~\citep{quijano2000coloniality,mohamed2020decolonial}. A leitura é post hoc,
de modo que reivindicamos a lacuna para a distinção Norte/Sul tal como
convencionalmente traçada, e não como afirmação geral sobre desenvolvimento.

\emph{O gradiente nulo é informativo, dentro de limites declarados.} Simulando
amostras de 25 países com correlações conhecidas, o desenho tem \npower\% de poder em
$\rho=0.55$ --- o limiar fixado de antemão --- e efeito mínimo detectável de
$\rho=\nmde$. Um gradiente da magnitude que nos comprometemos a chamar de
substancial está excluído; um fraco não está (poder de \npowerweak\% em $\rho=0.40$). Os dez
países acrescentados post hoc caem \emph{dentro} da faixa de IDH que os quinze
originais já cobriam, de modo que a extensão aumenta a densidade de observações, e
não a amplitude do preditor.

\subsection{Síntese}
\label{sec:res:sintese}

A Tabela~\ref{tab:sintese-hip} mapeia cada hipótese ao seu achado, ao mecanismo
candidato e a um estatuto de evidência explícito, de modo que nenhuma afirmação
seja lida como mais forte do que os dados sustentam. Três efeitos são sustentados e
sobrevivem a todo ataque de robustez da Seção~\ref{sec:res:robustez}; duas
expectativas pré-especificadas não se cumprem; e o mecanismo de corpus é
identificado dentro dos países apenas até o par cobertura-ou-desenvolvimento,
que este desenho não consegue separar.

\begin{table}[!ht]
\centering
\caption{Hipóteses, achados, mecanismos candidatos e estatuto de evidência.
Agrupada pelo que a evidência sustenta, não por número de hipótese.}
\label{tab:sintese-hip}
\footnotesize
\setlength{\tabcolsep}{8pt}
\renewcommand{\arraystretch}{1.15}
\renewcommand{\arraystretch}{1.15}
\begin{tabular}{@{}>{\raggedright\arraybackslash}p{0.235\linewidth}%
                  >{\raggedright\arraybackslash}p{0.315\linewidth}%
                  >{\raggedright\arraybackslash}p{0.235\linewidth}@{}}
\toprule
Hipótese & Achado & Mecanismo \\
\midrule
\multicolumn{3}{@{}l}{\emph{Sustentadas}} \\
\addlinespace[2pt]
H2 idioma           & $\nnativasigned$~pp; hindi $\nnativahi$     & corpus do idioma \\
H3 modelo regional  & $\delta=\nregdelta$, pior de 14   & escala sobre ajuste \\
H1 lacuna de camada & $\ntiergap$~pp; RC \nort{} em T1     & camada \\
\midrule
\multicolumn{3}{@{}l}{\emph{Sustentada, mecanismo não separável}} \\
\addlinespace[2pt]
H4 dentro do país   & $\beta=\nhfourbeta$ sozinha; $\nhfourjcob$ ($p=\nhfourjcobp$) com IDH & cobertura ou desenvolvimento \\
\midrule
\multicolumn{3}{@{}l}{\emph{Não sustentadas}} \\
\addlinespace[2pt]
H1 gradiente        & $\rho=\nrho$, abaixo de $0.55$ & desenvolvimento \\
H4 pré-especificada & proxy de corpus do idioma nulo & corpus do idioma \\
H6 persona          & DiD $\ndid$~pp, $p=\npersonap$      & --- \\
\midrule
\multicolumn{3}{@{}l}{\emph{Apenas descritiva}} \\
\addlinespace[2pt]
H5 aberto/fechado   & $\nhfive$~pp para fechados     & tipo de acesso \\
\bottomrule
\end{tabular}

\vspace{4pt}
\begin{minipage}{0.75\linewidth}
\footnotesize\emph{Nota.} O mecanismo de H2 repousa sobre três idiomas nativos e é
descritivo. A lacuna de camada de H1 é concentrada nas tarefas cuja resposta
depende de um fato nacional. H4 dentro do país descarta o canal do corpus da
língua, mas não separa cobertura de desenvolvimento, e não é identificação causal.
\end{minipage}
\end{table}
